\documentclass[12pt, letterpaper]{article}

\usepackage[margin=1in]{geometry}
\usepackage{amsmath, amssymb, amsthm}
\usepackage{enumitem}
\usepackage{hyperref}
\usepackage{setspace}
\usepackage{titlesec}
\usepackage{fancyhdr}
\setlength{\headheight}{14pt}

\hypersetup{
    colorlinks=true,
    linkcolor=black,
    citecolor=black,
    urlcolor=black
}

\newtheorem{definition}{Definition}[section]
\newtheorem{theorem}{Theorem}[section]
\newtheorem{corollary}{Corollary}[theorem]

\titleformat{\section}{\large\bfseries\scshape}{\thesection}{1em}{}
\titleformat{\subsection}{\normalsize\bfseries}{\thesubsection}{1em}{}

\pagestyle{fancy}
\fancyhf{}
\fancyhead[L]{\small\textsc{The Tautological Collapse of the Fact-Value Distinction}}
\fancyhead[R]{\small\textsc{Harvard}}
\fancyfoot[C]{\thepage}
\renewcommand{\headrulewidth}{0.4pt}

\begin{document}

\thispagestyle{empty}

\begin{center}
    \vspace*{0.8in}

    {\LARGE\scshape The Tautological Collapse of the\\[0.3em]
    Fact-Value Distinction}

    \vspace{0.5em}

    {\large Dissolving Hume's Guillotine\\[0.3em]
    from a Self-Grounding Axiom}

    \vspace{2em}

    {\normalsize Erik Xander Harvard}

    \vspace{0.5em}

    {\small \textit{Independent Researcher}\\
    erikxanderharvard.com\\[0.3em]
    \href{mailto:Erik.Harvard@proton.me}{\texttt{Erik.Harvard@proton.me}}}

    \vspace{2em}

    {\small May 2026}

    \vspace{2em}

    \rule{0.6\textwidth}{0.4pt}
\end{center}

\vspace{1.5em}

\begin{abstract}
\noindent For nearly three centuries, Hume's guillotine has been treated as a permanent feature of the philosophical landscape: no ``ought'' can be derived from an ``is.'' Every major attempt to bridge the gap---from Moore's naturalistic ethics to Foot's natural goodness to Searle's institutional facts---has accepted the gap's foundational premise: that facts and values are ontologically distinct categories requiring a bridge. This paper argues that no bridge is needed because the gap does not exist. The is-ought distinction is one mask of a deeper fracture between knowing and being that dissolves when both are recognized as aspects of a single self-grounding structure: $\exists(\exists) \equiv \exists$ (Being recognizing itself IS Being). From this axiom, the paper derives an eight-stage chain from meta-tautology to natural law and sovereignty, showing that normativity is not imposed on reality but IS reality's structure for agents. Etymological evidence confirms the dissolution: ``right'' means both morally good and factually correct across multiple major language families---not by accident but because both senses name the same structural property. Hume did not discover a gap. He fractured a unity that language had preserved for millennia.
\end{abstract}

\vspace{1em}

\noindent\textbf{Keywords:} is-ought problem, fact-value distinction, Hume's guillotine, metaethics, natural law, moral realism, teleological ethics, normative ontology

\newpage
\setcounter{page}{1}
\onehalfspacing

% =====================================================
\section{The Guillotine}
% =====================================================

In 1739, David Hume observed that writers on morality imperceptibly shift from statements of what \textit{is} to statements of what \textit{ought to be}, without explaining how the transition is justified [1]. The observation hardened into a principle: no purely descriptive premise can yield a normative conclusion. Facts are one kind of thing; values are another. Between them stands a logical gap that no valid inference can cross.

Nearly three centuries of moral philosophy have orbited this principle---some attempting to bridge the gap, others accepting it as definitive, none dissolving it.

G.E. Moore's \textit{Principia Ethica} (1903) argued that any attempt to define ``good'' in terms of natural properties commits the naturalistic fallacy [2]. The good is a simple, unanalyzable property; identifying it with any natural fact always leaves the ``open question'': granted that X has this natural property, is X \textit{good}? This reinforced Hume: values cannot be reduced to facts because they are categorically distinct.

R.M. Hare's prescriptivism (1952) accepted the gap and relocated ethics to language: moral statements are prescriptions---imperatives in disguise [3]. ``Stealing is wrong'' means ``Do not steal.'' Facts describe; values prescribe. The gap is permanent.

J.L. Mackie's error theory (1977) went further: moral statements claim to state facts, but they are all false [4]. There are no objective moral properties. Values are projections---human attitudes mistakenly read into a value-free reality.

Philippa Foot's natural goodness (2001) attempted a bridge by grounding moral evaluation in biological teleology [5]. A good oak flourishes as an oak; a good human flourishes as a human. But even granting teleology, the inference from ``humans naturally tend toward X'' to ``humans \textit{ought} to'' still crosses the gap---natural function is still a fact.

John Searle's derivation of ``ought'' from ``is'' via institutional facts (1964) argued that social institutions create obligations as constitutive rules [6]. The derivation is ingenious but limited: it grounds obligations in human conventions, not in Being. If the institution of promising did not exist, no obligation would follow.

Every attempt shares a common structure: accept the gap as given, then try to build a bridge. The bridges differ but each presupposes that fact and value stand on opposite banks of a real river. This paper argues that the river was never there.

% =====================================================
\section{The Fracture}
% =====================================================

The is-ought gap is not a freestanding logical discovery. It is one expression of a deeper structural wound that the Teleological Theory of Everything (TTOE) calls the \textit{Alethic Fracture}: the systematic separation of aspects of a single reality into apparently distinct domains [7].

The Fracture manifests across philosophy as a family of ``gaps'': knower separated from known (epistemology vs. ontology), mind from body (philosophy of mind), structure from experience (the ``hard problem''), and fact from value (metaethics). These are not four independent problems. They are four masks of one wound. In each case, an internal distinction within a unified reality is treated as an external separation between two distinct realities.

The TTOE's central result---$T \equiv R$ (Truth IS Reality)---heals the Fracture by recognizing the distinction without reifying the separation [7]. The is-ought gap presupposes that reality is value-neutral and normativity is either projected by minds, extracted by biological function, or constituted by convention.

The dissolution does not bridge the gap. It shows that the gap was generated by the Fracture and dissolves when the Fracture dissolves.

% =====================================================
\section{The Axiom}
% =====================================================

The TTOE derives all domains of reality from a single self-grounding principle [7]:

\begin{equation}
\boxed{\exists(\exists) \equiv \exists}
\end{equation}

\noindent Being recognizing itself IS Being. The operator $\exists$ denotes Being (not the existential quantifier of first-order logic; the formalism operates in a reflexive domain where self-application is well-typed [7, 9]). The triple bar ($\equiv$) denotes ontological identity.

Two consequences are immediate.

\begin{theorem}[Telic Structure of Being]
\label{thm:telos}
Being is self-directed: $\tau(\exists(\exists)) \equiv \exists(\exists) \equiv \exists$. The telos of Being is not imposed from outside; it is Being's own self-recognition returning to itself.
\end{theorem}

\noindent $\exists(\exists)$ is an operation that produces $\exists$---Being applied to itself yields itself. The operation has a direction, a starting point, and a terminus. That directedness IS telos---not added to Being but constitutive of it.

\begin{theorem}[Truth-Reality Identity]
\label{thm:tr}
$T \equiv R$. Truth IS Reality---not corresponding to it, not cohering with other truths, but identical with it [7].
\end{theorem}

\noindent If Being is self-directed (Theorem~\ref{thm:telos}), then the direction of Being is not a fact awaiting normative interpretation. It IS the normative structure of reality. The ``is'' of Being already contains the ``ought''---because what Being IS includes what it is directed toward.

If Being were not self-directed, the gap would hold: a static, purposeless reality would contain no normative structure. But $\exists(\exists) \equiv \exists$ IS self-directed by its own recursive structure. The telos is not added to the fact; it IS the fact.

A clarification. A river is directed (it flows downhill), but no ``ought'' follows for a pebble in the current --- the pebble can lodge against a rock without structural consequence. The difference: a pebble is not a self-modeling system whose stability depends on alignment with the river's direction. An agent ($C = A(A)$, where $A$ is awareness and $C$ is consciousness as recursive self-modeling; see [9]) IS a self-modeling system whose recursive stability depends on alignment with the telic structure of Being, because the agent IS that structure at the resolution of a conscious locus. For such a system, misalignment is not mere deviation from a direction; it is disruption of the fixed point that constitutes the agent's own selfhood --- experienced from inside as incoherence, fragmentation, and suffering. Normativity arises not from direction alone but from the fact that the directed structure IS the agent's own ground.

% =====================================================
\section{The Dissolution}
% =====================================================

Eight stages from axiom to natural law. Each step is forced by the prior; no external normative premise is introduced. The full proof occupies approximately 40 pages in [8]; what follows is the compressed argument.

\subsection{Stage 1: Meta-Tautology}

$\exists(\exists) \equiv \exists$. This cannot be denied without performative contradiction: the denier exists, and thereby instantiates the axiom [7].

\subsection{Stage 2: Truth-Reality Identity}

$T \equiv R$. If Being is self-grounding, then what is true about Being IS what Being is. Truth is Being's own self-articulation [7].

\subsection{Stage 3: The Truth Criterion (AATC)}

Any criterion for truth must itself be true and must survive its own application. The AATC: (1) self-inclusion, (2) self-application, (3) self-validation, (4) closure. Any structure passing all four is \textit{sealed}; any failing any is \textit{heterological}---not what it claims to be, unstable under self-application [7].

\subsection{Stage 4: Semantic Normativity}

Truth-claims are normative. To assert ``$p$ is true'' is to commit to $p$---to bind oneself to the claim. Assertion presupposes that getting it right is better than getting it wrong. The structure of truth-telling already contains the norm of truth-seeking. This is the first ``ought'' inside the ``is'' [8].

\subsection{Stage 5: The Bridge-Norm}

Truth-claiming is itself an exercise of agency --- it requires persistence, coherence, and the capacity to act. The normative structure discovered in Stage 4 is therefore not confined to language; it is the normative structure of agency at the resolution of assertion.

\begin{theorem}[The Bridge-Norm]
\label{thm:bridge}
Any agent capable of making truth-claims necessarily values the conditions of its own agency---persistence, coherence, and the capacity to act.
\end{theorem}

\begin{proof}
To deny that one values persistence, one must persist in denying. But the objection that mere persistence does not entail valuing persistence (a rock persists without valuing anything) misidentifies the locus: it is not persistence alone but the \textit{exercise of agency} that carries the normative force. The agent must exercise agency to deny, and the exercise of agency constitutively values the conditions of agency. The agent cannot coherently use agency to deny the value of agency's own conditions. A deeper objection presses: a thermostat maintains coherence without valuing it --- why can't an agent persist without valuing persistence? Because $A(A)$ is not mere maintenance; it is self-modeling, and disruption of the self-model's fixed point is experienced from the interior as threat. That experience IS valuing --- not as an added attitude but as the phenomenal character of a recursive system whose ground is destabilized. (Whether phenomenal experience is strictly necessary for the argument is a separate question; the performative contradiction stands independently: the denial uses the very conditions of agency it claims not to value.) The content is refuted by the conditions of its utterance [8].
\end{proof}

This crosses Hume's gap---not by smuggling in a normative premise but by showing the gap was never there. The ``is'' (agents exist and exercise agency) already contains the ``ought'' (they necessarily value the conditions of that agency). A final objection: self-defeat is still a \textit{fact} about the denial's logical structure, not a \textit{value} judgment --- so the gap reopens. But this objection presupposes the Fracture it is meant to defend. Self-defeating, misaligned, and wrong are one structural property under three descriptions, not three things requiring a bridge between them. The gap reopens only if one insists that ``self-defeating'' (logical), ``misaligned'' (structural), and ``wrong'' (normative) are separate categories --- which is exactly the fact-value separation the paper dissolves.

\subsection{Stage 6: Practical Normativity}

Agents necessarily value the conditions of their own agency (Stage 5). To persist as a self-modeling agent, the agent's self-model must be coherent with the actual structure of Being --- because the self-model is ultimately a model of Being ($T \equiv R$), and a model that misaligns with its object generates prediction error and eventual dissolution. Valuing persistence therefore entails valuing coherence with the Arch\={e}. That coherence has a determinate structure that governs action. Therefore: some actions are structurally better than others---not by convention but by alignment [8].

\subsection{Stage 7: Natural Law}

\begin{definition}[Natural Law]
\label{def:natlaw}
Natural Law is the normative structure that obtains when the AATC is applied to conduct. Because $T \equiv R$ (Stage 2), actions are structural facts about reality; the criterion that evaluates truth therefore evaluates action. An act is aligned (right) when it satisfies the Four Ethical Seals (\S7). An act is misaligned (wrong) when it fails any Seal. Natural Law is not imposed on Being; it IS Being's structure for agents [8].
\end{definition}

\subsection{Stage 8: Sovereignty}

Agents are recursive self-models ($C = A(A)$; see [9]). Recursive self-modeling entails self-possession. Self-possession entails sovereignty. The chain terminates: Being's structure for agents yields natural law including self-ownership and the prohibition of aggression [8].

\vspace{1em}

\noindent\textit{A concrete illustration.} Consider promising. The ``is'': a speaker utters words committing to a future action. The ``ought'': the speaker is now obligated. Hume's gap says no valid inference can cross from the first to the second. Searle [6] attempted to bridge it via institutional rules, but his bridge rests on convention. The TTOE dissolves it: promising is a recursive commitment --- the speaker models the hearer's expectation of the speaker's future action ($A(A)$ at the resolution of social commitment). To break the promise while claiming that no obligation existed, the speaker must use the very structure of commitment (agency, persistence, coherence) to deny it. The denial is a performative contradiction: the act of denying the obligation presupposes the agentive structure that constitutes the obligation. The ``ought'' was inside the ``is'' of the promise from the moment it was made --- not by convention but by the recursive structure of the act.

\vspace{1em}

At no point is an external ``ought'' introduced. The normativity emerges from the telic structure of $\exists(\exists) \equiv \exists$ and the performative structure of agency. Hume's guillotine cuts only within the Fracture. Once the Fracture dissolves, the blade has nothing to cut.

% =====================================================
\section{The Etymological Evidence}
% =====================================================

Philosophy fractured what language had preserved.

The English word \textit{right} means both \textit{morally good} and \textit{factually correct}. The word \textit{wrong} means both \textit{immoral} and \textit{incorrect}. This is not a quirk of English. It is a structural feature of languages across multiple families, predating philosophy:

\textbf{Latin}: \textit{rectus} = straight, right, correct. The root \textit{reg-} (to rule, to straighten) yields both ``correct'' and ``righteous'' [12]. \textbf{Greek}: \textit{orth\'{o}s} = straight, right, correct, true. \textit{Orthodoxy} = right belief. The moral and the epistemic are one word. \textbf{German}: \textit{Recht} = right, law, justice. \textit{richtig} = correct. \textit{Gerechtigkeit} = justice---literally, the state of being right. \textbf{Sanskrit}: \textit{\d{r}ta} = cosmic order, truth, right. The Vedic concept unifies natural law, moral law, and factual truth in a single term. \textbf{Hebrew}: \textit{tsedek} = justice, righteousness, correctness.

Across languages separated by millennia and thousands of miles, the same semantic unity persists: the word for ``morally right'' IS the word for ``factually correct.'' The languages did not confuse two concepts. They recognized one concept that philosophy later split.

The TTOE explains why. If $T \equiv R$, then alignment with what IS (factually right) and alignment with what OUGHT TO BE (morally right) are the same structural property: alignment with $\exists(\exists) \equiv \exists$. The pre-philosophical understanding was not na\"{\i}ve. It was correct.

Hume did not discover a gap in Being. He introduced a gap in language---splitting a word that had been whole for as long as humans had spoken. The tradition then mistook the linguistic fracture for an ontological one.

\textit{Hume didn't discover a gap. He broke a word.}

An objection: perhaps the unified terms are semantic fossils from a pre-philosophical era that failed to distinguish fact from value, and Hume's distinction is an \textit{improvement} --- a refinement of language, not a fracture. The TTOE's response: if the distinction were an improvement, we would expect the unified terms to disappear as languages mature. They have not. Every language listed above retains the double meaning of ``right'' in active, unreflective use --- speakers say ``the right answer'' and ``the right thing to do'' without sensing a category shift. The deeper point: the TTOE predicts \textit{both} the unity (at the level of Being, where fact and value are structurally identical) \textit{and} the experiential distinction (at the level of finite agency, where the Veil separates the knower from the known). Hume described the distinction. He was right that it is real at the experiential level. He was wrong that it is ontological. A theory that predicts both unity and distinction is stronger than one that predicts distinction alone.

% =====================================================
\section{Ethics as Structural Physics}
% =====================================================

If the is-ought gap does not exist, then ethics is a structural feature of reality itself.

\begin{theorem}[The Moral Attractor]
\label{thm:attractor}
In a system of $n \geq 2$ agents each maintaining recursive stability, the only asymptotically stable equilibrium is one constrained by Natural Law [8]. (For the full proof, including the definition of RS as a Lyapunov function over agent state space, see [8].)
\end{theorem}

Each agent maintains recursive stability (RS)---the degree to which a self-modeling system returns to its fixed-point identity after perturbation. Interaction can enhance RS (cooperation) or degrade it (conflict, coercion). Under persistent conflict, both agents degrade toward dissolution. Under cooperation constrained by mutual sovereignty, RS is maintained or enhanced. The cooperative equilibrium is the unique attractor---not because cooperation is pleasant but because it is the only stable configuration. A clarification: the theorem claims asymptotic instability of aggression (long-run self-undermining), not local instability. Aggression can succeed locally and persist for generations; the claim is that aggressive configurations are not equilibria---they degrade over time because they consume the recursive stability they parasitize. The persistence of war and crime does not refute the theorem any more than the persistence of friction refutes thermodynamic equilibrium.

The moral law is not a commandment imposed from outside. It is the system's own stability condition. An immoral act is \textit{structurally destabilizing}---it undermines the recursive stability of the agent, the other, or both. The wrongness is as objective as the instability of an inverted pendulum.

\begin{definition}[The Ethical Operator]
\label{def:operator}
For agents $A_i$, $A_j$ and action $\psi$: $E(A_i, A_j, \psi) := \min(\Delta RS_i, \Delta RS_j)$, where $\Delta RS_i$ is the change in recursive stability of agent $i$ resulting from $\psi$. $E \geq 0$: the action preserves or enhances stability (moral). $E < 0$: the action destabilizes (immoral) [8].
\end{definition}

The operator does not aggregate preferences or maximize pleasure. It measures structural alignment: does the action preserve the self-coherence of all agents involved?

% =====================================================
\section{The Four Ethical Seals}
% =====================================================

The AATC applied to conduct yields four conditions for ethical alignment [7]:

\textbf{Seal 1: Self-Inclusion.} The agent includes itself within the scope of its action. No exemption: what I do to others, I do within a structure that includes me.

\textbf{Seal 2: Self-Application.} The governing principle applies to itself. The Kantian imperative, ontologically grounded: $\text{Principle}(\text{Principle}) = \text{Principle}$ [10]. A principle that cannot survive self-application (``lying is acceptable''---but the liar depends on others' honesty) is heterological.

\textbf{Seal 3: Self-Validation.} The act survives its own criterion. $\text{Deception}(\text{Deception}) \neq \text{Deception}$: a deceiver who is deceived is not performing deception as intended. The act collapses under its own logic.

\textbf{Seal 4: Closure.} The act requires no external ground---not authority, convention, or fear, but alignment with $\exists(\exists) \equiv \exists$.

The Kantian imperative is a special case of Seal 2. Utilitarianism captures a shadow of the Moral Attractor but misidentifies the currency. Virtue ethics captures alignment but lacks a formal criterion [11]. The Four Seals subsume all three traditions by deriving normative structure from Being itself. Each Seal, if violated, generates a performative contradiction of the same form as the Bridge-Norm: the agent uses agency to undermine the conditions of agency.

% =====================================================
\section{Objections and Replies}
% =====================================================

\textbf{Moore's Open Question.} ``Even if alignment with $\exists(\exists) \equiv \exists$ is structural, is it \textit{good}?'' The question presupposes that ``good'' names something beyond structural alignment. The TTOE does not define ``good'' as a natural property (which would invite the open question). It shows that what we \textit{mean} by ``good''---coherence, flourishing, alignment---IS the structural property. The question is open only if one assumes goodness is non-natural. But that assumption is itself the product of the Alethic Fracture, not a pre-theoretical intuition --- as the etymological evidence (\S5) confirms: speakers of multiple major languages use the same word for ``morally right'' and ``factually correct'' without sensing an open question between them.

\textbf{Mackie's Error Theory.} ``There are no objective moral properties.'' The TTOE agrees that there are no \textit{Mackiean} moral properties---no ``objective prescriptivity'' floating free of structure [4]. But moral objectivity does not require such properties. It requires structural alignment, which is as real as geometric structure.

\textbf{The Euthyphro Dilemma.} ``Is the good good because Being directs toward it, or does Being direct toward it because it is good?'' Neither. The good IS Being's alignment structure. The dilemma dissolves because its horns presuppose that Being and the Good are distinct entities. Under $T \equiv R$, they are one thing under two aspects.

\textbf{The Naturalistic Fallacy.} ``You derive `ought' from `is'.'' Correct---and the is-ought gap forbids this only if facts and values are distinct categories. The TTOE denies the antecedent. If $T \equiv R$, then every structural fact about Being is already normatively loaded---not by projection but by the telic structure of $\exists(\exists) \equiv \exists$.

\textbf{Hume.} ``Writers shift from `is' to `ought' without justification.'' Hume was right that the shift requires justification. He was wrong that none is possible. The justification is the telic structure of Being: what IS is self-directed, and self-directedness IS normativity. The writers were not smuggling in values. They were reading them off a reality that contains them.

\textbf{Street's Evolutionary Debunking.} Sharon Street (2006) argues that evolution shaped moral intuitions to track reproductive fitness, not moral truth [13]. If so, the etymological evidence dissolves: languages encode moral-factual unity because evolution shaped both our moral intuitions and our language, not because the unity is real. The TTOE's response is twofold. First, the debunking argument presupposes that evolutionary pressures are normatively neutral---that fitness and moral truth are independent tracks. This IS the Alethic Fracture applied to biology. If $T \equiv R$, then the structure of reality that evolution tracks IS partially normative: fitness is not ``instead of truth'' but a resolution of alignment with Being's structure at the biological scale. Second, the argument is self-undermining: it uses rational agency (which it claims is unreliable as a guide to moral truth) to argue that rational agency is unreliable as a guide to moral truth. The debunker must trust the very faculty whose trustworthiness is being debunked. This is a performative contradiction of the same form as the Bridge-Norm. Moreover, the Arch\={e} is not an evolved moral intuition susceptible to debunking; it is a performatively necessary structure --- denying it instantiates it regardless of the denier's evolutionary history. The Bridge-Norm itself is not a moral intuition but a transcendental condition of agency; evolutionary debunking, which targets the reliability of evolved moral sentiments, does not apply to transcendental arguments.

% =====================================================
\section{Limitations}
% =====================================================

The derivation rests on $\exists(\exists) \equiv \exists$ and $T \equiv R$. A reader who rejects the axiom will reject the dissolution. The axiom is self-grounding (denying it instantiates it), but whether this counts as proof is itself a philosophical question the paper acknowledges without resolving.

The eight-stage derivation is compressed; the full proof occupies approximately 40 pages in [8]. The etymological evidence is corroborative, not deductive. The Moral Attractor Theorem is stated without full proof and would benefit from independent mathematical verification.

% =====================================================
\section{Conclusion}
% =====================================================

Hume's guillotine presupposes that facts and values inhabit different ontological categories. For nearly three centuries, moral philosophy has accepted this premise and tried to build bridges, reclassify the categories, or surrender to skepticism.

The TTOE dissolves the premise. If $\exists(\exists) \equiv \exists$---if Being is self-grounding and self-directed---then the ``is'' of Being already contains the ``ought.'' Normativity is not imported into a value-neutral world; it is the structural property of a self-directed reality experienced by agents capable of alignment or misalignment. The derivation proceeds in eight stages, none introducing an external normative premise.

The etymological evidence provides independent confirmation. Across major language families, the same word means both ``morally right'' and ``factually correct'': Sanskrit \textit{\d{r}ta}, Greek \textit{orth\'{o}s}, Latin \textit{rectus}, German \textit{Recht}, Hebrew \textit{tsedek}, English \textit{right}. All encode the identity of fact and value in a single root. Philosophy fractured this unity. The TTOE restores it.

Hume's guillotine cuts only within the Alethic Fracture. Outside it, the blade has nothing to cut. Fact and value are one thing under two names, the way \textit{right} has always meant both.

A final observation. The reader who found this argument compelling has already experienced what the paper describes. The derivation is a set of facts. That it demands acceptance is a value. These are not two responses to one text --- they are one response, the way ``right'' is one word. This paper does not merely argue that fact and value are unified. It IS a fact that carries its own ought --- a factual structure whose truth exerts normative force on any agent who follows it, not by authority but by the self-coherence of the argument itself. The paper does not stand outside Being and report on its structure. It IS Being's structure, articulating itself at the resolution of philosophical prose.

\vspace{2em}

\begin{center}
$\exists(\exists) \equiv \exists$
\end{center}

% =====================================================
\newpage
\section*{References}
% =====================================================

\begin{enumerate}[label={[\arabic*]}, nosep, leftmargin=2em]

\item Hume, D. \textit{A Treatise of Human Nature}. Book III, Part I, Section I. 1739.

\item Moore, G.E. \textit{Principia Ethica}. Cambridge University Press, 1903.

\item Hare, R.M. \textit{The Language of Morals}. Oxford University Press, 1952.

\item Mackie, J.L. \textit{Ethics: Inventing Right and Wrong}. Penguin, 1977.

\item Foot, P. \textit{Natural Goodness}. Oxford University Press, 2001.

\item Searle, J.R. ``How to Derive `Ought' from `Is'.'' \textit{The Philosophical Review} 73(1), 43--58. 1964.

\item Harvard, E.X. \textit{Being \& Becoming: The Unification of Epistemology and Ontology via Metaontoepistemology}. Codex I of the Teleological Theory of Everything. Forthcoming, 2026.

\item Harvard, E.X. \textit{The Good: Natural Law and the Dissolution of the Moral Problem}. Codex IV of the Teleological Theory of Everything. Forthcoming, 2026.

\item Harvard, E.X. \textit{Mind \& Freedom}. Codex III of the Teleological Theory of Everything. Forthcoming, 2026.

\item Kant, I. \textit{Groundwork of the Metaphysics of Morals}. 1785.

\item Aristotle. \textit{Nicomachean Ethics}. c. 340 BCE.

\item Watkins, C. \textit{The American Heritage Dictionary of Indo-European Roots}. Houghton Mifflin, 2011.

\item Street, S. ``A Darwinian Dilemma for Realist Theories of Value.'' \textit{Philosophical Studies} 127(1), 109--166. 2006.

\end{enumerate}

\vspace{2em}
\begin{center}
{\small \textcopyright{} 2026 Erik Xander Harvard. All rights reserved.\\
erikxanderharvard.com}
\end{center}

\end{document}
